% ---------------------------------------------------------------------
\chapter{Zusammenarbeit und Konflikte}
\label{kap:zusammenarbeit}

\begin{lernziele}
\begin{itemize}[nosep]
  \item mit mehreren Geräten oder Personen am selben Repository arbeiten,
  \item verstehen, wie ein Konflikt entsteht,
  \item einen Konflikt Schritt für Schritt auflösen,
  \item Konflikte durch Arbeitsgewohnheiten vermeiden.
\end{itemize}
\end{lernziele}

\section{Mehrere Geräte, ein Repository}

PC und Laptop besitzen je einen eigenen Klon. Sie wissen nichts voneinander; verbunden sind
sie nur über GitHub:

\begin{center}
\begin{tikzpicture}[
    geraet/.style={draw=hsluBlau, fill=hsluBlau!6, rounded corners=3pt, thick,
                   minimum width=2.6cm, minimum height=1cm, font=\small\bfseries},
    weg/.style={-{Stealth[length=2.5mm]}, thick},
    lab/.style={font=\ttfamily\footnotesize, fill=white, inner sep=1pt}
  ]
  \node[geraet] (pc) at (0,0)    {PC};
  \node[geraet] (gh) at (5,0)    {GitHub};
  \node[geraet] (lt) at (10,0)   {Laptop};
  \draw[weg] ([yshift=2.5mm]pc.east)  -- node[lab, above=1mm] {push} ([yshift=2.5mm]gh.west);
  \draw[weg] ([yshift=-2.5mm]gh.west) -- node[lab, below=1mm] {pull} ([yshift=-2.5mm]pc.east);
  \draw[weg] ([yshift=2.5mm]gh.east)  -- node[lab, above=1mm] {pull} ([yshift=2.5mm]lt.west);
  \draw[weg] ([yshift=-2.5mm]lt.west) -- node[lab, below=1mm] {push} ([yshift=-2.5mm]gh.east);
\end{tikzpicture}
\end{center}

Mehrere Personen funktionieren genauso, nur dass auf jedem Gerät jemand anderes arbeitet.
Daraus folgt eine einfache Regel:

\begin{merke}
\textbf{Zu Beginn \cmd{git pull}, am Ende committen, pushen, mergen.} Wer das einhält, hat
auf allen Geräten denselben Stand und erlebt Konflikte nur noch selten.
\end{merke}

\subsection*{Was Git überträgt und was nicht}

\begin{itemize}
  \item \textbf{Mit Git übertragen} wird alles, was im Repository committet ist, auch
        Einstellungen wie \cmd{.vscode/settings.json} und Skripte wie \cmd{setup/einrichten.ps1}.
  \item \textbf{Nicht übertragen} werden installierte Programme, Anmeldungen wie die von
        \cmd{gh} und alles, was in \cmd{.gitignore} steht.
\end{itemize}

Darum enthält dieses Repository ein Einrichtungs-Skript. Auf jedem neuen Gerät genügt nach
dem Klonen ein Befehl:

\begin{lstlisting}[style=befehle]
powershell -ExecutionPolicy Bypass -File setup\einrichten.ps1
\end{lstlisting}

\begin{teile}
  powershell              & startet eine neue PowerShell, die das Skript ausführt \\
  -ExecutionPolicy Bypass & erlaubt Skripte nur für diesen Aufruf; Windows blockiert sie sonst standardmässig \\
  -File                   & die auszuführende Skript-Datei \\
\end{teile}

\begin{praxis}
Dieses Vorgehen ist in Firmen Standard und heisst oft \emph{Onboarding-Skript}: Wer neu
ins Team kommt, klont das Repository, startet ein Skript und kann am selben Tag
produktiv arbeiten. Was ein Gerät braucht, steht als Code im Repository statt in den
Köpfen einzelner Personen.
\end{praxis}

\section{Wie ein Konflikt entsteht}

Git führt Änderungen meist selbstständig zusammen, auch wenn zwei Personen dieselbe Datei
bearbeitet haben, solange es \emph{verschiedene Stellen} sind. Ein \begriff{Konflikt}
entsteht nur, wenn \textbf{dieselbe Stelle} auf zwei Wegen unterschiedlich geändert wurde.
Git kann dann nicht wissen, welche Fassung richtig ist, und fragt den Menschen.

Typisches Beispiel mit zwei Geräten:

\begin{enumerate}
  \item Am PC wird in \cmd{README.md} Zeile 5 geändert und gepusht.
  \item Am Laptop wurde vorher \emph{kein} \cmd{git pull} gemacht. Dort wird dieselbe Zeile
        anders geändert und committet.
  \item \cmd{git pull} am Laptop meldet: \cmd{CONFLICT (content): Merge conflict in README.md}.
\end{enumerate}

\section{Einen Konflikt auflösen}

Git markiert die betroffene Stelle direkt in der Datei:

\begin{lstlisting}[style=ausgabe]
<<<<<<< HEAD
Wöchentlich zwei Übungen lösen.
=======
Täglich eine Übung lösen.
>>>>>>> 7f02c1a (README: Lernrhythmus angepasst)
\end{lstlisting}

\begin{teile}
  <<<<<<< HEAD    & Beginn der eigenen Fassung, also des aktuellen Stands (hier: Laptop) \\
  =======         & Trennlinie zwischen den beiden Fassungen \\
  >>>>>>> 7f02c1a & Ende der anderen Fassung; dahinter Hash und Nachricht des Commits, der von GitHub kam \\
\end{teile}

Vorgehen in vier Schritten:

\begin{enumerate}
  \item Datei öffnen und die Stelle suchen. VS Code hebt sie farbig hervor und bietet
        Schaltflächen wie \emph{Accept Current} oder \emph{Accept Incoming}.
  \item Die gewünschte Endfassung herstellen: eine der beiden behalten oder beide
        sinnvoll kombinieren.
  \item \textbf{Alle Markierungszeilen} (\cmd{<<<<<<<}, \cmd{=======}, \cmd{>>>>>>>}) entfernen.
  \item Den aufgelösten Stand speichern und abschliessen:
\end{enumerate}

\begin{lstlisting}[style=befehle]
git add README.md                      # "Konflikt ist gelöst" melden
git commit -m "Konflikt in README aufgelöst"
git push
\end{lstlisting}

\begin{fehler}
Markierungszeilen werden übersehen und mitcommittet. Danach steht \cmd{<<<<<<< HEAD} im Code,
und das Programm startet nicht mehr. Vor dem \cmd{git add} darum die Datei nach
\cmd{<<<<<<<} durchsuchen (Strg+F).
\end{fehler}

\begin{erfahrung}
Ein Konflikt ist kein Fehler und nichts Gefährliches. Git hat nichts verloren, es fragt
lediglich nach. Wer unsicher ist, kann den Vorgang jederzeit mit \cmd{git merge --abort}
abbrechen; danach ist alles wie vor dem \cmd{pull}.
\end{erfahrung}

\section{Konflikte vermeiden}

\begin{itemize}
  \item Vor Arbeitsbeginn \cmd{git pull}.
  \item Häufig pushen, statt tagelang lokal zu sammeln.
  \item Kurzlebige Branches: je kürzer ein Branch lebt, desto weniger weicht er von \cmd{main} ab.
  \item Im Team absprechen, wer an welchen Dateien arbeitet.
\end{itemize}

\begin{zusammenfassung}
\begin{itemize}[nosep]
  \item Geräte und Personen sind nur über GitHub verbunden: am Anfang \cmd{pull}, am Ende \cmd{push}.
  \item Ein Konflikt entsteht nur, wenn dieselbe Stelle unterschiedlich geändert wurde.
  \item Auflösen: Stelle bereinigen, Markierungen löschen, \cmd{add}, \cmd{commit}, \cmd{push}.
  \item \cmd{git merge --abort} bricht ab und stellt den vorherigen Zustand wieder her.
\end{itemize}
\end{zusammenfassung}
